% Section II: Related Work -- fingerprinting with nearest neighbors, the
% dataset, tree models, neural networks on UJIIndoorLoc, leakage and grouped
% CV, and the gap. What each paper did or reported is taken from
% docs/literature_review.md and docs/literature_matrix.csv (source locations
% are given there). No result of our own appears in this section.
\section{Related Work}\label{sec:related}

\emph{Fingerprinting.} Surveys~\cite{liu2007survey,he2016wifi} cover the
field. The classic RADAR system~\cite{bahl2000radar} located a scan by the
recorded scans nearest to it in signal space, the nearest-neighbor
rule~\cite{cover1967nearest} behind kNN. Horus~\cite{youssef2005horus} modeled
the variation of the signals with probabilities. Both were tested on the
authors' own small testbeds (median error 2.94~m for RADAR, below 0.6~m on
average for Horus), so their errors cannot be compared with other studies.

\emph{The dataset.} UJIIndoorLoc~\cite{torres2014ujiindoorloc} was presented
as the first public Wi-Fi fingerprinting database with several buildings and
floors. Its validation file was recorded months after the training file,
partly by new users and phones, which makes it a realistic
test~\cite{torres2015comprehensive}. Torres-Sospedra
et~al.~\cite{torres2015comprehensive} compared many kNN settings on it, but
chose the best one on the file they report, so that result is optimistic.

\emph{Tree models.} A comparison of classifiers on UJIIndoorLoc found kNN the
most suitable, with boosting and bagging bringing a decision tree close to
it~\cite{bozkurt2015comparative}. We test two tree ensembles, random
forests~\cite{breiman2001random} and XGBoost~\cite{chen2016xgboost}, tuned and
with our own features.

\emph{Neural networks.} Nowicki and Wietrzykowski~\cite{nowicki2017loweffort}
classified building and floor on UJIIndoorLoc with a deep neural network. Kim
et~al.~\cite{kim2018scalable} trained one network with building, floor and
location outputs. CNNLoc~\cite{song2019cnnloc} used a convolutional neural
network, but its headline numbers are the best values of different runs, not
of one model.

\emph{Grouped CV.} For grouped data, random CV can seriously underestimate the
error, and grouped (block) CV is the better
choice~\cite{roberts2017crossvalidation}. UJIIndoorLoc shows this too: a
network that got building and floor right for 98.2\% of a random split of the
training file did so for only 91.1\% of the validation
file~\cite{nowicki2017loweffort}.

\emph{Gap.} We found no study that compares tuned kNN, random forests and
XGBoost on UJIIndoorLoc with every choice made by user-group CV and the
validation file used only once, as the test set. Phase~1 fills this gap. Our
models still predict building, floor and position separately; multi-task
learning~\cite{caruana1997multitask}, which learns related tasks together
through a shared representation, is the subject of Phase~2.
